\section{Comparison pages with the operations retained}
\label{p:chapter}

The relation rectangle gives one proof for every pointed Barr-exact
Mal'cev category. There is another useful way to watch the comparison
proceed: let simple factors persist through a sequence of pages until
maps pair them. We can carry this calculation out simultaneously for
groups and modules, without pretending that arbitrary group maps behave
like linear maps.

For this chapter a group with operators means a group equipped with a
fixed family of endomorphisms. Subgroups and normal subgroups under
discussion must be invariant under all those endomorphisms, and maps
commute with them. Modules are included by their additive groups and the
operators $x\mapsto rx$ for every scalar $r$. Quotient simplicity is
simplicity with this structure retained. Ordinary groups are the case
of an empty family of extra operators. This is a concrete common instance
of the preceding framework; the page construction below is asserted for
this instance, with composition series as input.

\subsection{The normality that makes the pages possible}

Write two composition series as
\[
1=A_0\nml A_1\nml\cdots\nml A_m=G,
\qquad
1=B_0\nml B_1\nml\cdots\nml B_n=G.
\]
We use multiplicative notation, including for the identity subgroup $1$.
Extend both filtrations constantly below zero and beyond their last
indices. Put $S_i=A_i/A_{i-1}$ and $T_j=B_j/B_{j-1}$ for the real
indices $1\leq i\leq m$, $1\leq j\leq n$.

\begin{lemma}\label{p:normality}
A subnormal subgroup of a simple group with operators is either trivial
or the whole group. For every integer $k$ and every real index $j$,
\[
 B_{j-1}(A_k\cap B_j)
 \quad\text{is either }B_{j-1}\text{ or }B_j.
\]
In particular this product is an invariant normal subgroup of $B_j$.
Similarly the image of $A_i\cap B_k$ in $S_i$ is either trivial or all
of $S_i$.
\end{lemma}
\begin{proof}
Take a finite normal chain from a proper subnormal subgroup $H$ to a
simple group $S$. The term immediately before the first occurrence of
$S$ is proper and normal in $S$, so it is trivial. The earlier terms,
including $H$, must be trivial. This proves the first claim with
invariance under the operators included.

Intersecting a normal chain with a fixed subgroup preserves adjacent
normality: conjugating an element of $U\cap K$ by an element of
$V\cap K$ stays in both $U$ and $K$ when $U\nml V$.
Taking images under a homomorphism also preserves adjacent normality,
since $f(v)f(u)f(v)^{-1}=f(vuv^{-1})$. Intersections and images here
remain invariant under every operator. Thus $A_k\cap B_j$ is subnormal
in $B_j$, by intersecting the tail of the $A$-chain with $B_j$; the
constant extensions at the endpoints cause no change to this assertion.
Its image in the simple factor $T_j$ is trivial or everything. The
preimage of that image is exactly $B_{j-1}(A_k\cap B_j)$, proving
the stated alternatives. The same argument with the roles exchanged
proves the last assertion.
\end{proof}

This check is essential. A nonzero homomorphism between simple groups
need not be onto: an inclusion $C_2\hookrightarrow A_5$ is an immediate
counterexample. The maps constructed below have subnormal, and in fact
normal, images for a specified reason. That is what makes the cancellation
argument valid.

\subsection{Terms, maps, and the next page}

For $r\geq1$ define two rows
\begin{align}
 U_r(i)&=
 \frac{A_i\cap B_{n+i-r}}{A_{i-1}\cap B_{n+i-r}},
       &&1\leq i\leq m,\label{p:source}\\
 V_r(j)&=
 \frac{B_j}{B_{j-1}(A_{j+r-n-1}\cap B_j)},
       &&1\leq j\leq n.\label{p:target}
\end{align}
The same formulas give trivial groups outside the real index ranges.
Every source denominator is invariant and normal in its numerator by
intersection. Every target denominator is invariant and normal by
Lemma~\ref{p:normality}. Consequently these are actual quotient groups
with all their operators, not just sets of cosets.

For $j=n+i-r$, inclusion followed by quotient defines
\begin{equation}\label{p:differential}
 d_{r,i}:U_r(i)\longrightarrow V_r(j),\qquad [x]\longmapsto[x].
\end{equation}
Its source denominator is $A_{i-1}\cap B_j$, and this is contained in
the target denominator $B_{j-1}(A_{i-1}\cap B_j)$. Thus the map is
well-defined, is a homomorphism, and preserves the operators. Each source
has one designated target; each target has one designated incoming source.

\begin{proposition}[The page transition]\label{p:transition}
With $j=n+i-r$, the inclusion and quotient maps give canonical
isomorphisms
\[
 U_{r+1}(i)\cong\ker d_{r,i},\qquad
 V_{r+1}(j)\cong V_r(j)/\im d_{r,i}.
\]
The image in the second formula is an invariant normal subgroup.
\end{proposition}
\begin{proof}
Put $H=A_i\cap B_j$, $D=A_{i-1}\cap B_j$, and $K=B_{j-1}$.
Thus the map in question is $H/D\to B_j/KD$. The subgroup $D$ is
normal in $H$, $K$ is normal in $B_j$, and both $KD$ and $KH$ are
normal in $B_j$ by Lemma~\ref{p:normality}. If $j$ is outside the real
target range, the constant extension gives $B_{j-1}=B_j$, so these
normality assertions hold automatically. This also covers transitions
whose designated source or target lies outside the initial rows.

We first compute its kernel. Since $D\leq H$,
\[
                     H\cap KD=(H\cap K)D.
\]
Indeed, if $h=kd\in H$ with $k\in K$ and $d\in D$, then
$k=hd^{-1}\in H\cap K$. The reverse inclusion follows by multiplying
elements already in $H$ and in $KD$. Restricting the quotient map to
$H\cap K$ therefore gives a surjection onto the kernel with kernel
$D\cap K$. Hence
\[
 \ker d_{r,i}\cong\frac{H\cap K}{D\cap K}
 =\frac{A_i\cap B_{j-1}}{A_{i-1}\cap B_{j-1}}
 =U_{r+1}(i).
\]
This isomorphism is induced by inclusion; all its operations are retained.

The image is $KH/KD$, a normal subgroup of $B_j/KD$. The next target
projection has precisely this kernel, so the compatible iterated quotient
map gives
\[
 \frac{V_r(j)}{\im d_{r,i}}
 \cong \frac{B_j}{KH}
 =\frac{B_j}{B_{j-1}(A_i\cap B_j)}
 =V_{r+1}(j).
\]
The last equality uses $i=j+r-n$. These are explicit maps between
quotients, so their inverses preserve the operators as well.
\end{proof}

The transition can be displayed as two short exact sequences:
\[
\begin{tikzcd}[column sep=small]
1\arrow[r]&U_{r+1}(i)\arrow[r,hook]&U_r(i)
 \arrow[r,two heads]&\im d_{r,i}\arrow[r]&1,\\
1\arrow[r]&\im d_{r,i}\arrow[r,hook]&V_r(j)
 \arrow[r,two heads]&V_{r+1}(j)\arrow[r]&1.
\end{tikzcd}
\]
The first row retains the source part not yet seen in the target. The
second removes the target part accounted for by that map. This is the
two-row page mechanism $E_{r+1}\cong H(E_r,d_r)$, with the group
quotients justified by the normality proof. We need no addition of maps
to perform this particular calculation.

\begin{theorem}[Every layer eventually finds its partner]\label{p:cancellation}
The pages above give a bijection between the simple factors $S_i$ and
$T_j$, together with isomorphisms preserving all the operators. The proof
uses neither Schreier refinement nor Jordan--H\"older as an input.
\end{theorem}
\begin{proof}
The source $U_r(i)$ is isomorphic to the image of $A_i\cap B_{n+i-r}$
in $S_i$, by the kernel of the restricted quotient map. By
Lemma~\ref{p:normality} it is trivial or isomorphic to $S_i$.
The target denominator in~\eqref{p:target} is either $B_{j-1}$ or
$B_j$, so $V_r(j)$ is trivial or isomorphic to $T_j$.

The image of $d_{r,i}$ is normal by Proposition~\ref{p:transition}.
A nontrivial map therefore surjects onto its simple target. Its kernel
is an invariant normal subgroup of its simple source; it cannot be the
whole source, so it is trivial. Consequently every nontrivial
differential is an isomorphism. The transition then removes precisely
that source and target. When the differential is trivial, the transition
retains both terms by its canonical isomorphisms.

On the first page, $n+i-1\geq n$ for every real source index, so
$U_1(i)=S_i$. For the targets $j-n\leq0$, giving $V_1(j)=T_j$.
If $G\neq1$, then $m,n\geq1$. On every page $r\geq m+n$ we have
$n+i-r\leq0$ and $j+r-n-1\geq m$ at real indices. Every term is
therefore trivial. Each initial factor must have disappeared exactly
once, by an isomorphism differential, and each such differential removes
one source and one target. Following the transition isomorphisms to that
page gives the required pairwise factor isomorphisms. For $G=1$, both
lists are empty and the conclusion is immediate.
\end{proof}

Notice the dependency: simplicity is part of the supplied filtrations;
invariance of their lengths is the conclusion. No page count assumes
that the two rows start with the same number of terms.

\subsection{A nonabelian crossing and the module specialization}

Let $G=A_5\times C_2$ and use the two series with intermediate subgroups
$A_5\times1$ and $1\times C_2$. Their first-page factors have opposite
orders. The first differential at $(i,j)=(1,2)$ identifies
$A_5\times1$ with $G/(1\times C_2)$. It removes the two $A_5$ factors.
The second page has no nontrivial differential, but the process has not
terminated. On the third page, at $(i,j)=(2,1)$, the actual differential
is the identity on $1\times C_2$, since both $U_3(2)$ and $V_3(1)$
are that group. Its source is identified with the original factor
$G/(A_5\times1)$ by the restricted quotient map
$1\times C_2\to G/(A_5\times1)$. Thus it pairs the two remaining
$C_2$ factors. The fourth page is trivial. A quiet page need not be the
last page.

For completeness, the simplicity of $A_5$ used only in this illustration
can be checked by its conjugacy classes. Its nonidentity elements are
the $15$ double transpositions, $20$ three-cycles, and $24$ five-cycles.
Their centralizers in $A_5$ have orders $4,3,5$, respectively: commuting
permutations permute cycles of equal length; on a three- or five-cycle they
are powers of that cycle, with the odd transposition on the two unused
letters excluded in the three-cycle case; the even permutations commuting
with a double transposition form its four-element Klein group.
Counting conjugates by cosets of the centralizer therefore gives class
sizes $15,20,12,12$. A normal subgroup is a union of these classes and
the identity, and its order divides $60$ because its cosets partition
$A_5$. None of the proper nontrivial sums of $1$ and a selection from
$15,20,12,12$ divides $60$. Hence no proper nontrivial normal subgroup
exists. This finite check is not a prerequisite of the page theorem.

For modules the products in~\eqref{p:target} are sums, and the first
row of~\eqref{p:source} is a subquotient of modules. Put
$F_pX=A_{p-n}$, $F_pY=B_p$, and take $f=\id_G:X\to Y$ on the
underlying module. The page formulas in the next chapter recover exactly
$U_r(i)$ at $p=n+i$ and $V_r(j)$ at $q=j$, with the differential
landing at $q=p-r$. The next chapter proves that calculus for an
\emph{arbitrary filtered module map}, where the pages may have nonsimple
terms and nonzero limiting kernel or cokernel. The present calculation
explains why its composition-series cancellation can also be carried
out for groups with operators.

The scope is precise. We proved the relation refinement theorem for all
pointed Barr-exact Mal'cev categories. We proved this explicit page
calculation for composition series of groups with operators. We have not
asserted that every filtered group complex supports the ordinary abelian
spectral-sequence machinery. General nonabelian exact-couple theories
require additional constructions; one primary reference is
\href{https://arxiv.org/pdf/1007.0632}{Grandis, \emph{Homotopy spectral
sequences}}, Sections~2--4.
